\section{Introducción}
Los nuevos estándares post-cuánticos del NIST, como ML-KEM~\cite{fips203}, se apoyan en
problemas de celosías, en particular en variantes de LWE. Cuando la dimensión es grande nadie
conoce una forma eficiente de resolverlos, ni siquiera con un ordenador cuántico. En este trabajo
hemos atacado instancias pequeñas y medianas de GGH~\cite{ggh}, LWE~\cite{regev} y
Ring-LWE~\cite{lpr,tema11} para ver por nosotros mismos qué los hace difíciles y cuándo dejan de
serlo.

Usamos la notación del enunciado. Los vectores son filas y una base $B$ genera la celosía
$L(B)=\{xB : x\in\mathbb Z^n\}$, todas las combinaciones enteras de sus filas. En SVP se busca el
vector no nulo más corto y en CVP el punto de la celosía más cercano a un objetivo $t$. En GGH el
criptograma $c=xB'+e$ es un punto de la celosía movido un poco por el error $e$, así que descifrar
es resolver un CVP. En LWE la clave pública es $b=As+e \pmod q$: sin el error sería un sistema de
ecuaciones normal, y es el error lo que lo complica.

Lo hemos programado todo en Python~3, con fracciones exactas y sin bibliotecas de celosías. Las
mismas herramientas (LLL y enumeración) nos sirven para GGH y para LWE. En Ring-LWE aprovechamos
una debilidad de los parámetros de las tareas que detectamos al analizar sus datos.
